\scanpage{14}

ens fermé de $I$ contenant $i$) \deleted{à l'image}\inserted{correspond} au plus petit ens. fermé \uline{saturé} de $X$ contenant $X_i$, c'est donc $\overline X_i$ par 1). Donc $j\in\overline{\{i\}}$ ssi $X_j\subset\overline X_i$ i.e. $\overline X_j\subset\overline X_i$ : la relation d'incidence est la relation de spécialisation. La condition \inserted{$2^\circ)$ : $\overline X_i=\overline X_j\Longrightarrow X_i=X_j$} i.e. $i=j$ \uline{signifie} donc que $I$ est ``présobre'' \deleted{\unsure{condition}} i.e. $\overline{\{i\}}=\overline{\{j\}}\Longrightarrow i=j$. Si $\varphi:X\to I$ est l'application de passage au quotient, la condition $1^\circ)$ \uline{signifie} que $J\mapsto\varphi^{-1}(J)$ commute à l'adhérence (ou encore, à l'intérieur) i.e. que $\varphi$ est une application \uline{ouverte}.

Ainsi, on voit que les décompositions cellulaires \inserted{(finies)} de $X$ correspondent aux \deleted{\unsure{esp}} espaces quotients \inserted{finis} de $X$ qui sont présobres, et tels que $X\to X/R$ soit \uline{ouverte}.

\S\ 3 \deleted{\unsure{Théorème|Admissibilité}} \inserted{\uline{Décompositions cellulaires \unsure{Fin-admiss.|$F$-admiss.} : th. d'existence}}

Si $F$ est comme \unsure{ds|dans} \S\ 1, une décomposition cellulaire de $X$ est dite $F$-admissible (ou partition $F$-admissible) si ses feuilles sont $\in F$, i.e. si $F_{\mathcal P}\subset F$.

\mspage{15}
